\section{Introduction}\label{sec:intro}

Access-to-information law assumes that public administrations know how they work. Freedom of
information regimes, open data mandates and machine-readable publication duties are all written
as though an accurate account of administrative practice sits somewhere in a filing system,
waiting to be released. \citet{brauneis2018algorithmic} tested that assumption directly, filing
42 open-records requests across 23 states, and found the binding obstacle was not trade secrecy
or exemption. It was that the records had never been created. There was nothing to disclose.

Administrative process systems make the gap partially solvable. A modern electronic process
system records every action taken on every case: which unit received it, what was done, when it
moved on. The trace exists. What does not exist is the account of what those traces mean in
aggregate: which kinds of procedure the administration actually runs, and how each of them
actually flows. Administrations do maintain a declared taxonomy of procedure types, but that
taxonomy is a classification applied at opening, by the person opening the case, and its
relationship to observed practice is an empirical question rather than a given.

In the corpus studied here, that relationship is weak in a specific and diagnosable way. A
single catch-all label covers 30.5\% of the concluded cases. Compare an emergent partition derived from the
event log against the declared taxonomy and the agreement is modest (ARI 0.202); remove that one
label and agreement rises to 0.352. The official taxonomy is not uninformative everywhere. Its
disagreement with practice is concentrated in a junk drawer, which is the characteristic failure
of a classification maintained for filing rather than for analysis.

This paper derives the missing record and shows how such a record has to be evaluated. From the
event log of a state-level secretariat of a Brazilian state administration, we build an event
knowledge graph of 5{,}243 concluded administrative cases, group them into emergent process
types using the text of each case's opening document together with its routing and action
history, and learn a typical path for each type. The result is a catalogue: a set of named
process types, each with a representative flow, derived from the administration's own records
and regenerable from them.

\subsection*{Contributions}

\textbf{An auditable as-is catalogue of process types, offered as a transparency instrument.}
The artefact is what the machine-readable publication duty in Brazil's access-to-information
law contemplates, and what \citeauthor{brauneis2018algorithmic}'s requesters could not obtain.
Two properties make it usable rather than merely produced. It is non-random: permuting cluster
membership while holding cluster count, size distribution and evaluated population fixed drops
downstream next-step accuracy from 0.726 to 0.494, so the grouping carries the signal rather
than the shape of the partition. And it is reproducible: every clustering that is defined as a function of the feature matrix refits exactly
from its stamped feature matrix, and every number reported in this paper is regenerated from
artefacts by a pipeline stage that fails if a generated table drifts, and the load-bearing
figures in the running text are held to the artefacts by a second check that fails in either
direction.

\textbf{An external evaluation protocol, with controls, for judging such a catalogue.} The
principle of scoring a clustering by the downstream model it enables is established
\citep{deweerdt2013actitrac}; we extend it from process-model fitness to next-activity
prediction, apply it externally to partitions produced independently by four algorithm families
across two trace alphabets, and add two controls this literature does not routinely report. The
controls matter more than the extension. Measured against the conventional within-cluster
baseline, per-cluster modelling appears to gain 0.083; measured against a single global model
trained on the pooled corpus, the gain is 0.013. Both numbers are correct measurements of
different quantities, and reporting only the first is what makes a modest result look strong.

The two contributions are deliberately independent. The catalogue claim rests on the shuffled
control and on reproducibility, not on the size of the downstream gain, so the modest gain does
not undercut it.

\subsection*{Research questions}

\begin{description}
  \item[RQ1] Can an auditable catalogue of emergent process types be derived from the event log
    and opening-document text, and does it capture structure the declared taxonomy misses?
  \item[RQ2] Does the choice of clustering algorithm matter to downstream utility, and do
    internal cluster validity indices predict it --- or do they measure something else?
  \item[RQ3] Do per-cluster typical paths generalise --- covering what real cases do without
    fabricating steps --- and do they improve on a global model, a random partition, and a real
    representative case?
  \item[RQ4] Does constrained decoding provide the compliance guarantee an audit artefact
    requires, and at what cost?
  \item[RQ5] Does anonymisation degrade the structure or its downstream utility?
\end{description}

Several answers are negative, and they are reported as findings. Algorithm choice separates the
clusterings very little: sixteen spanning four algorithm families fall within 0.0079 of each
other once scored on the cases all of them evaluate, though the interval around that spread is
too wide for an equivalence claim. The internal validity indices that would ordinarily choose
between them measure something else: silhouette tracks how much of the corpus a partition
retains, and has no association with accuracy once cluster size is held fixed. A real
representative case remains a strong baseline that the generated paths do not dominate. These results constrain what the method can claim, and they
are the reason the positive results are credible.

Section~\ref{sec:related} situates the work; Section~\ref{sec:legal} sets out the normative
framing; Sections~\ref{sec:method} and \ref{sec:setup} describe the corpus, the method and the
experimental design; Section~\ref{sec:results} reports results by research question;
Sections~\ref{sec:discussion}--\ref{sec:conclusion} discuss, bound and conclude.
